

\documentclass[a4paper,aps,superscriptaddress,floatfix,nofootinbib,longbibliography, onecolumn]{revtex4-1}

% ----- Packages ------
\usepackage{amsmath, amsthm, amssymb} % Standard math packages
\usepackage{algpseudocode} % Pseudocode package
\usepackage{comment} % Comment environment
\usepackage{graphicx}% Include figure files
\usepackage{dcolumn}% Align table columns on decimal point
\usepackage{bm}% bold math
\usepackage{hyperref}% add hypertext capabilities
\usepackage[mathlines]{lineno}% Enable numbering of text and display math
%\linenumbers\relax % Commence numbering lines
\usepackage{float}
%\usepackage[showframe,%Uncomment any one of the following lines to test 
%scale=0.7, marginratio={1:1, 2:3}, ignoreall,% default settings
%text={7in,10in},centering,
%margin=1.5in,
%total={6.5in,8.75in}, top=1.2in, left=0.9in, includefoot,
%height=10in,a5paper,hmargin={3cm,0.8in},
%]{geometry}



\usepackage{xcolor}
\newcommand{\change}[1]{{#1}}
%\newcommand{\newchange}[1]{{\color{red}{#1}}}
\newcommand{\newchange}[1]{{#1}}
  


\usepackage{comment}


\renewcommand{\theequation}{S\arabic{equation}}
\setcounter{equation}{0}
\renewcommand{\thefigure}{S\arabic{figure}}
\setcounter{figure}{0}
\renewcommand{\thetable}{S\arabic{table}}
\setcounter{table}{0}


% ----- Packages ------
\usepackage{amsmath, amsthm, amssymb} % Standard math packages
\usepackage{algpseudocode} % Pseudocode package
\usepackage{comment} % Comment environment
\usepackage{graphicx}% Include figure files
\usepackage{dcolumn}% Align table columns on decimal point
\usepackage{bm}% bold math
\usepackage{hyperref}% add hypertext capabilities
\usepackage[mathlines]{lineno}% Enable numbering of text and display math
%\linenumbers\relax % Commence numbering lines
\usepackage{float}




\usepackage{xcolor}
%\newcommand{\change}[1]{{#1}}
%\newcommand{\newchange}[1]{{\color{red}{#1}}}
%\newcommand{\newchange}[1]{{#1}}

\begin{document}

\title{Modeling plant disease spread via high-resolution human mobility networks \\
Supplementary Materials}

\author{Varun K. Rao}
\affiliation{Center for Complex Networks and Systems Research, Luddy School of Informatics, Computing, and Engineering, Indiana University, Bloomington, Indiana 47408, USA}
\affiliation{School of Public Health, Indiana University, Bloomington, Indiana 47408, USA}
\author{Ryan Higgs}
\affiliation{Onside, Christchurch,  New Zealand.}
\author{Hautahi Kingi}
\affiliation{Onside, Christchurch,  New Zealand.}
\author{Filippo Radicchi}
\affiliation{Center for Complex Networks and Systems Research, Luddy School of Informatics, Computing, and Engineering, Indiana University, Bloomington, Indiana 47408, USA}
\author{Santo Fortunato}
\affiliation{Center for Complex Networks and Systems Research, Luddy School of Informatics, Computing, and Engineering, Indiana University, Bloomington, Indiana 47408, USA}
\author{Maria Litvinova}
\email{malitv@iu.edu}
\affiliation{School of Public Health, Indiana University, Bloomington, Indiana 47408, USA}
%TC:incbib

\maketitle

\section{Numerical Integration of the Mean-Field Approximation}

Denote with $I_i(t)$ the fraction of hectares that are infected in farm $i$ at time $t$, and with 
$S_i(t)$ the fraction of hectares that are not infected in farm $i$ at time $t$. We clearly have
$I_i(t) + S_i(t) = 1$ for all $i= 1, \ldots, N$ and all $t \geq 0$.

The farm-based mean-field equations for the SI dynamics are
\begin{equation}
\begin{array}{ll}
\frac{d \, I_i(t)}{dt} = & \beta_{\mathrm{w}} \, I_i(t) S_{i}(t) \, \Theta \left[I_i(t) - I_i^{\min} \right] +
\beta_{\mathrm{b}} \, S_{i}(t) \,  \sum_{j=1}^N w_{j \to i}(t)  \, I_{j}(t) \, \Theta \left[I_j(t) - I_j^{\min} \right] 
\end{array}
\; ,
%\frac{d \, I_i(t)}{dt} = \beta_{w} \, I_i(t) S_{i}(t) + \beta_{b} \, S_{i}(t) \,  \sum_{k=1}^N w_{k \to i}(t,Z)  I_{k}(t) \; 
\label{eq:cont}
\end{equation}
where $\beta_{w}$ and $\beta_{b}$ are the transmission rates within and between farms, respectively.  $\Theta(x)$ is the Heavyside function defined such that $\Theta(x) = 1$ if $x > 0$ and 
$\Theta(x) = 0$ otherwise. To be considered infected and contribute to the spreading dynamics, property $i$ should have a level of infection that is above the threshold  $I_{i}^{\min} = \left(H_i \, 625 \, \mathrm{plants \, hect}^{-1}\right)^{-1}$, where $625$ is a realistic estimate of the average number of plants per hectare~\cite{NZKGI2024}. 
$w_{j \to i}$ is proportional to the number movements between properties $j$ and $i$ observed at time $t$.
%$\omega^{(m)}_{j \to i}$ quantifies the average number of daily movements between properties $j$ and $i$ observed in month $m$.
%$\beta$s  are the "free" parameters that should be calibrated based on the available data (from literature or an ongoing outbreak) and will absorb the differences in the absolute values between $I_i(t)$ and $\sum_{k=1}^N w_{k \to i}  I_{k}(t)$.
%$\beta_{b}$ should always be small enough so that $( \beta_{w} \, I_i(t) + \beta_{b} \,  \sum_{k=1}^N w_{k \to i}  I_{k}(t) ) < 1 $ for all $i$.

We can also write the above as a finite-difference equation
\begin{equation}
I_{i}(t+dt) = I_{i}(t) + \beta_{w} \, dt \, \, I_i(t)  S_{i}(t)\Theta \left[I_i(t) - I_i^{\min} \right] + \beta_{b} \, dt \, \, S_{i}(t) \,  \sum_{j=1}^N w_{j \to i}(t)  \, I_{j}(t) \, \Theta \left[I_j(t) - I_j^{\min} \right] 
\label{eq:fd}
\end{equation}
%Now, 
%we numerically integrate
To perform the numerical integration, we require that  

\begin{equation}
    \beta_{w} \, dt \, \, I_i(t)  S_{i}(t)\Theta \left[I_i(t) - I_i^{\min} \right]  +  \beta_{b} \, dt \, \, S_{i}(t) \,  \sum_{j=1}^N w_{j \to i}(t)  \, I_{j}(t) \, \Theta \left[I_j(t) - I_j^{\min} \right]  \leq 1 - I_i(t)\; .
\end{equation}

%The condition is on the time increment $dt$ that can not be arbitrarily large, rather
leading to the constraint for the time increment:
\begin{equation}
    dt  \leq \frac{1 - I_{i}(t)}{\beta_{w} I_i(t)  S_{i}(t) \Theta \left[I_i(t) - I_i^{\min} \right] + \beta_{b}  \, \, S_{i}(t) \,   \,  \sum_{j=1}^N w_{j \to i}(t)  \, I_{j}(t) \, \Theta \left[I_j(t) - I_j^{\min} \right]} \; .
\end{equation}
Such a condition means that Eq.~(\ref{eq:fd}) can effectively be used to numerically integrate Eq.~(\ref{eq:cont}).
Since the above must be true for all farms, we need to impose

\begin{equation}
     dt  \leq \min \, \bigg( T_{\mathrm{inc}} ,\min_{i=1, \ldots, N} \, \frac{1 - I_{i}(t)}{\beta_{w} I_i(t)  S_{i}(t) + \beta_{b}  \, \, S_{i}(t) \,  \sum_{k=1}^N \frac{w_{k \to i}}{\tau}  I_{k}(t)} \; \bigg).
\end{equation}
     
where $T_{\mathrm{inc}} = 1$ to generate at least one data point per day.
%is the minimum daily time increment, so we set $T_{inc} = 1.$


\section{Infection Data}
The number of monthly cumulative infections and incidence from the 2010 Psa-V outbreak in New Zealand was digitized from Greer and Saunders \cite{greer2012costs} and is listed in Table \ref{tab:incidence_ma_simplified}. Raw data values are shown as well as the values for the trailing 2-month moving average. Using this incidence data we calibrated our model.
\begin{table}[H]
\centering
\caption{Monthly raw counts and two-month moving average (MA(2)) incidence and cumulative infections.}
\label{tab:incidence_ma_simplified}
\renewcommand{\arraystretch}{1.2}
\setlength{\tabcolsep}{6pt}
\begin{tabular}{|c|cc|cc|}
\hline
\textbf{Month} 
& \multicolumn{2}{c|}{\textbf{Raw Counts}} 
& \multicolumn{2}{c|}{\textbf{2 Month Moving Average}} \\
\cline{2-5}
 & Incidence & Cumulative 
 & Incidence & Cumulative \\
\hline
Feb.\ 2011 & 50  & 50   & 50  & 50   \\
Mar.\ 2011 & 25  & 75   & 25  & 75   \\
Apr.\ 2011 & 22  & 97   & 23.5  & 98.5  \\
May.\ 2011 & 16  & 113  & 19  & 117.5  \\
Jun.\ 2011 & 10  & 123  & 13  & 130.5  \\
Jul.\ 2011 & 15  & 138  & 12.5  & 143  \\
Aug.\ 2011 & 86  & 224  & 50.5  & 193.5  \\
Sep.\ 2011 & 88  & 312  & 87  & 280.5  \\
Oct.\ 2011 & 111 & 423  & 99.5 & 380  \\
Nov.\ 2011 & 467 & 890  & 289 & 669  \\
Dec.\ 2011 & 48  & 938  & 257.5 & 926.5  \\
Jan.\ 2012 & 31  & 969  & 39.5  & 966  \\
Feb.\ 2012 & 101 & 1070 & 66  & 1032 \\
Mar.\ 2012 & 71  & 1141 & 86  & 1118 \\
\hline
\end{tabular}
\end{table}
\section{Gridsearch Calibration}
    The model contains four free parameters which require calibration: the within-property spreading rate $\beta_w$, the between-property spreading rate $\beta_b$, the detection threshold $D$, and the initial seed property $s_0$. We test $1,300$ different seed properties located in the Bay of Plenty region and test several different values of the detection threshold $D$, where set $D = \{ 0.2,0.4,0.6,0.8\}$.
    
    To find an appropriate range of values for $\beta_w$ and $\beta_b$,  we first conducted a coarse exploratory gridsearch using logarithmic increments and searched over $[10^{-3},10^0]$. For each pair of  dynamical parameters, $(\beta_w,\beta_b)$,  the average RMSE was calculated using all values of $D$ and $s_0$. The RMSE landscape shown in Figure \ref{fig:log-hmap} indicates that the best parameter configurations are between $[10^{-2}, 10^{-1}]$ for all combinations of $\beta_w,\beta_b$, taking into account different data types and temporal aggregations. Informed by this finding, we performed a finer-grained gridsearch varying $\beta_w,\beta_b$ (Figure ~\ref{fig:lin-hmap}  over the range $[10^{-2}, 10^{-1}]$ using increments of $0.01$ across the set values of $D$ and $s_0$ . 

    \begin{equation}
         |x_c - y_c(\theta)| < \max \{ \epsilon, z* x_c \}
         \label{eq:supp-filter}
    \end{equation}
    
    The average RMSE values shown in both gridsearches (Fig ~\ref{fig:log-hmap} and Fig ~\ref{fig:lin-hmap}) were calculated across the set of configurations $G$, consisting of all possible combinations of $\beta_b,\beta_w,D,s_0$ that were discussed for each gridsearch.  Given $G$, and each configuration entry $\theta$  for all free parameters, we  filtered  configurations using the heuristic defined in the main text to prevent overfitting to the incidence curves.
    
    
    We slightly redefine the filtering heuristic in Equation \ref{eq:supp-filter}, using $\epsilon$ and $z$ as the two main variables that need to be set when performing the filtering. $\epsilon$ is utilized as a tolerance value while $z$ is the size of our uncertainty band. In the main text,  $\epsilon = 30$ and $z = 0.8$. A gridsearch was conducted to find the best values for $\epsilon$ and $x$  using the raw data in Figure ~\ref{fig:find-eps}. $\epsilon$ was varied from $10^{1}$ - $10^{2}$ in increments of $10^{1}$ while $z$ was varied from $10^{-1} - 1$ in increments of $10^{-1}$. The percentage of viable configurations that existed by fulfilling the conditions in Equation \ref{eq:supp-filter} was recorded for each combination. From that gridsearch, $\epsilon=30$ and $z=0.8$ are the first two values that produce the smallest percentage of viable configuration among all aggregations. See the main text for details on how we selected the best configuration for each aggregation type. 

    \begin{figure}[H]
    \centering
    \includegraphics[width=.8\linewidth]{figs/SM-PDF/log_hmap.pdf}
    \caption{\textbf{Logarithmic gridsearch calibration of $\beta_b,\beta_w$} We present the average RMSE for the various combinations of $(\beta_w,\beta_b)$ where we vary the values for these dynamical parameters from $10^{-3}$ to $10^{0}$ in logarithmic increments.  Each square in the heatmap is the expected RMSE value for the joint distribution of $\beta_w$ and $\beta_b$ over all the possible $s_0$ and $D$ combinations.  The black square within each heatmap represent the values for $\beta_b,\beta_w$ that were searched for in more depth in the linear gridsearch used in the main text. The RMSE values in the first column of heatmaps is calculated with reference to the raw data, while the values in the second column are calculated with respect to the two month moving average (MA(2)) data. In panel (A) we plot the RMSE for the yearly aggregation using the real data , in panel (B) we plot the average RMSE for the yearly network aggregation using the MA(2) data, in panel (C) we plot the average RMSE for the seasonal network aggregation using the raw data. while in panel (D) we plot the average RMSE for the seasonal aggregation using the MA(2) data. In panel(D) and (F) we plot the average RMSE for the monthly aggregation using the raw incidence data and MA(2) incidence data respectively .}
    \label{fig:log-hmap}
\end{figure}

% need to show linear gridsearch

\begin{figure}[H]
    \centering
    \includegraphics[width=.8\linewidth]{figs/SM-PDF/lin_hmap.pdf}
    \caption{\textbf{Linear gridsearch calibration of $\beta_b,\beta_w$} Similar to Figure ~\ref{fig:log-hmap} we plot the average RMSE for combinations of of $(\beta_w,\beta_b)$, where RMSE values in the first columns of the heatmaps are calculated with reference to the raw data, while the values in the second column are calculated with reference to the MA(2) data. We vary the values for these dynamical parameters from $10^{-2}$ to $10^{-1}$ respectively, in increments of $10^{-1}$.   }
    \label{fig:lin-hmap}
\end{figure}
    
     
    
    
    %The set of filtered configurations for each aggregation type $R$ is referred to the viable set $V_{R}$. 
    
   

    
    \begin{figure}[H]
        \centering
        \includegraphics[width=0.5\linewidth]{figs/SM-PDF/eps_vary.pdf}
        \caption{\textbf{Parameter gridsearch for $\epsilon$ and $z$}. The parameter values for $\epsilon$ and $z$ are varied here. We show the results of three different gridsearches, for all aggregations, using the raw data for the incidence curve to determine the . The darker the square in the heatmap, the higher percentage of viable configurations exist for that specified combination of $\epsilon$ and $z$. Panel (A) reports the percentage of viable configurations for the yearly aggregation, (B) and (C) are the same as (A) except for the seasonal and monthly aggregations respectively. }
        \label{fig:find-eps}
    \end{figure}

    
    Using the viable configurations for each aggregation and data type, we plot the posterior distribution of the parameters in Figure \ref{fig:post-dist}. We observe that there are strongly preferred values for all parameters and aggregations, barring the monthly aggregation for the raw data. For $\beta_w$ and $\beta_b$, specific values are clearly preferred. For the monthly raw aggregation, there is a bimodal distribution. 

\begin{figure}[H]
    \centering
    \includegraphics[width=\linewidth]{figs/SM-PDF/post_dist.pdf}
    \caption{\textbf{Posterior Distribution of Linear Gridsearch} Using the RMSE values generated from the linear gridsearch, seen in Fig. ~\ref{fig:lin-hmap}, we plot the posterior distribution for each parameter (except the seed, due to the number of seeds). These distributions are generated after creating the viable configuration set $V$ for each aggregation and data type (Raw and MA(2) incidence data) and limiting the viable configurations to the configurations that have the smallest $1\%$ of RMSE values among all configurations $V$. The best-fitting configurations are found from the configurations in $V$ that minimize the RMSE best for each configuration and dat type. The first columns RMSE values are calculated with reference to the raw data, while the second columns RMSE values are calculated with reference to the MA(2) data. Three different colors are used in each panel to reference different network aggregations. Blue represents the yearly aggregation, orange represents the seasonal aggregation, and green represents the monthly aggregation. Each row contains the posterior distribution for a different parameter. Panels (A) and (B) plot the distribution for $\beta_w$, panels (C) and (D) plot the distribution for $\beta_b$, and panels (E) and (F) plot the distribution for $D$. The RMSE values in panels (A), (C) and (E) are computed with respect to the raw data while in panels (B),(D), and (E) they are calculated using the MA(2) data. }
    \label{fig:post-dist}
\end{figure}

\section{Spatial Spread}
The complementary cumulative distribution function (CCDF) of movement length (i.e., distance between origin and destination property) is calculated and shown in Fig.~\ref{fig:ccdf-year}.  
$99.1$\% of movements occur between properties that are between $0-100$ km from each other,  with $71.2$ \% of movements occurring having a distance of $10$ km. From such a movement distribution, we can conclude that most movements in horticulture movement network occur between properties that are relatively close to one another.  This movement structure leads to the results in Figure~\ref{fig:prox-spread} where newly infected properties are infected within $10$ km of another previously infected property. 
While most movements are short-distance, long distance spread from the seed still occurs, as we find in Figure ~\ref{fig:max-seed}, where by the end of the outbreak the epidemic has spread more than $1,000$ km from where it started. 


% FIGURE: CCDF 
 \begin{figure}[H]
     \centering
     \includegraphics[width=.6\linewidth]{figs/SM-PDF/log_ccdf.pdf}
     \caption{\textbf{Complementary Cumulative Distribution Function (CCDF) of Movement Length.} }
     \label{fig:ccdf-year}
 \end{figure}


\begin{figure}[H]
    \centering
    \includegraphics[width=.6\linewidth]{figs/SM-PDF/prox_bars.pdf}
    \caption{\textbf{Distance of Newly Infected Property to the Closest Infected Property} Once a property becomes infected, we measure the distance in kilometers from that property to the closest infected property. An infected property $i$ is defined as one where the level of infection on it exceeds $D$ where $I_i > D$. We classify this distance using three categories: $0 - 10$ km, $11-20$ km, and $ > 20$ km. Most newly infected properties are within $0-10$ km of a previously infected property, and only $2.7 \%$ of properties are $ > 20$ km away from an infected property. }
    \label{fig:prox-spread}
\end{figure}

\begin{figure}[H]
    \centering
    \includegraphics[width=0.9\linewidth]{figs/SM-PDF/max_seed.pdf}
    \caption{\textbf{Maximum Spread Distance from Seed} Using the calibrated seasonal model, we plot the largest geodesic distance an infected property has from the initially infected seed property $s_0$. We define this distance as the spread distance. The node with the largest distance away from the seed could persist for months, as we observe that for consecutive months at a time the spread distance does not change in value. As reference, we plot the maximum and average movement distance between any two properties as well. By the end of the outbreak, the distance from the seed is very close to the maximum existing movement distance. }
    \label{fig:max-seed}
\end{figure}

\section{Wind Connectivity Matrix}
Leveraging wind reanalysis data spanning the years 2010-2012 from the ERA5 \cite{hersbach2023era5} dataset we utilize the R package windscape \cite{kling2021global} to generate a wind connectivity matrix for New Zealand. The ERA5 dataset compiles hourly high-resolution data over a latitude-longitude grid with a resolution of roughly 31 km. Using this matrix, we can estimate the wind travel time in hours between any two locations in New Zealand. First, we plot the wind-interconnectivity matrix in Figure \ref{fig:wind-map} over a map of New Zealand, where we show a subset of nodes from our yearly mobility network. We selected a node $i$ (in red) with the largest number of connections $j$ (nodes with at least one movement between $i$ and $j$, in white). The majority of these connections are within 20 hours of travel time from $i$ (assuming favorable conditions), but a few exceed that value. The wind conditions do not inhibit spreading; however, the observed patterns would make long-distance dispersal difficult in this example, as it would require more than 60 hours of pathogen survival and consistent favorable conditions. For privacy reasons, the latitude and longitude values are not shown. 

Using the wind-connectivity matrix, which we partially graph in Figure \ref{fig:wind-map},  we plot the relationship between wind travel time and distance between all $i$ in the yearly network and all of their connections in Figure \ref{fig:wind-dist}. We observe a strong, positive linear relationship between the two in Figure \ref{fig:wind-dist}, showing that as the distance between two properties grows, the travel time required for the transmission increases. These results show that while it requires less than a day for a pathogen to be transmitted through human mobility connections, it may require more than 10 days for the wind-based dispersal to reach the same location. This highlights the higher importance of incorporating human mobility networks for better understanding human-mediated long-distance dispersal.

The CCDF (Complementary Cumulative Distribution Function) in Figure \ref{fig:wind-ccdf} shows that when the distance between a property and its connections exceeds 20 km, there is a higher probability of the required travel time to exceed one day. These two figures highlight how adding dispersal via wind would lead to more spread within properties and between close properties.

\begin{figure}[H]
    \centering
    \includegraphics[width=0.9\linewidth]{revise_figs/wind_inbound_travel_time_zoom.png}
    \caption{\textbf{Windmap Example} Wind connectivity matrix between a highly-connected network node and all of its connections. The raster map color shows the amount of time for a wind particle to travel. The node is shown in red, its connections are shown in white.  }
    \label{fig:wind-map}
\end{figure}


\begin{figure}[H]
    \centering
    \includegraphics[width=0.8\linewidth]{revise_figs/wind_vs_distance.pdf}
    \caption{\textbf{Relationship between Wind Travel Time and Distance} Every node has a set of connections. This figure illustrates the relationship of the physical geodesic distance between a node and its connections and the wind travel time potentially needed for a wind-based dispersal to occur.}
    \label{fig:wind-dist}
\end{figure}



\begin{figure}[H]
    \centering
    \includegraphics[width=0.75\linewidth]{revise_figs/ccdf_wind.pdf}
    \caption{\textbf{CCDF of Wind Travel} The CCDF shows the wind travel time by distance class. Red dots refer to a property's connections within 10km, green dots refer to a property's connections 11-20km from it, while blue refers to a property's connections further than 20km from it. The y-axis denotes the probability of observing a value of the wind travel time $t$ greater than the listed y-value $T$. }
    \label{fig:wind-ccdf}
\end{figure}



\section{Comparing the Best Fits Based on Temporal Network Aggregations}

\begin{figure}[H]
    \centering
    \includegraphics[width=\linewidth]{revise_figs/best_fit_SM.pdf}
    \caption{Model fit for different network aggregations. Each panel shows the best-fitting configuration of each aggregation in a solid color. The lines in lighter/transparent color represent the other possible viable configurations with the same values for $\beta_b,\beta_w, D$ as the best-fitting configuration but varying seed locations $s_0$. The first row of figures shows fits for the raw data, while the second row shows fits for the moving average data. Panels (A) and (D) show yearly aggregations for the raw and moving average data, respectively. Panels (B) and (E) show seasonal aggregation for the raw and moving average data, respectively. Panels (C) and (F) show monthly aggregation for the raw and moving average data, respectively. }
    \label{fig:SM-viable}
\end{figure}

In Figure 4 of the main text, we show the best-fitting configurations along with uncertainty bands generated based on viable parameter configurations. For each aggregation, we selected a best-fitting parameter configuration for $\beta_w,\beta_b, D$ based on our calibration procedure, and simulated the outbreak for each $s_0$. Therefore, the uncertainty band represents the variation of the trajectories due to the selection of different seed properties. 

In Figure \ref{fig:SM-viable}, we show the individual trajectories of these simulations in lighter colors and the best-fit parameter selection in darker colors for each network aggregation, providing a more detailed variation of Figure 4. This figure also highlights the differences in the model fit for the raw data and the moving average. For the moving average dataset, we observe how the graph better represents the peak as the network aggregation resolution changes from yearly to seasonal to monthly. 

\section{Varying Initial Infection Ratio}
We calibrate the model assuming that the seed property's ($s_0$) infection level ($I_{s_0}$) is set to $0.1$. Figure \ref{fig:vary-a} shows the temporal dynamics of the outbreak when the initial amount we infect the seed with, $I_{s_0}$, is varied. We denote $I_{s_0}$ as $\alpha$ for this experiment. We vary $\alpha$ using the calibrated model based on the seasonal network. When calibrating the model $\alpha = 0.1$ , and find that increasing $\alpha$ increases the size of the outbreak while not having a large effect on the incidence dynamics. This result highlights the small effect that $\alpha$ has on outbreak size, and how it does not have a large effect on outbreak dynamics. 
  
    \begin{figure}[H]
        \centering
        \includegraphics[width=0.6\linewidth]{figs/SM-PDF/vary_alpha.pdf}
        \caption{\textbf{Dependence of Model's Outcome on Initial Seed Value} We vary the initial seed value $\alpha$  and plot the effect this has on the incidence curve in panel (A) and on the cumulative total infections shown in panel (B). We set $\alpha = 0.1$ when calibrating the model. Here we vary $\alpha$ from 0.1 to 0.9 in increments of 0.2 in both panels. In panel (A) we plot the real incidence curve as a reference, and in panel (B) the real cumulative infection data is plotted. For all graphs, the seasonal network is used as well as the corresponding parameters of best fit}
        \label{fig:vary-a}
    \end{figure}



%\printbibliography

\bibliography{refs.bib} 


\end{document}